\documentclass[11pt]{article}
\usepackage{amsmath,amssymb,amsthm}
\newtheorem{lemma}{Lemma}[section]
\newtheorem{corollary}[lemma]{Corollary}
\newtheorem{definition}[lemma]{Definition}
\newtheorem{remark}[lemma]{Remark}
\newtheorem{theorem}[lemma]{Theorem}
\newcommand{\WR}{W_R}
\newcommand{\WI}{W_I}
\newcommand{\ch}{\operatorname{ch}}
\newcommand{\KA}{A}
\newcommand{\CH}{\operatorname{CH}}
\newcommand{\Cl}{\operatorname{cl}}

\title{GST Birth Strike: the Component-Charge Law at the Vertex\\
(the boss's component-birth laws executed as construction engine) --- Round 17}
\author{GST Hodge campaign (SKILL V6.1 MODE B, ULTRA; the trajectory plan's Strike 1)}
\date{2026-10-27}

\begin{document}
\maketitle

\section*{Scope box (read first)}

\textbf{Claimed here:} Lemma B1 (the W-dead total law: every NS-tuned
codimension-2 complete-intersection total on the discriminator has Weil charge
zero --- the charge channel is only the signed component distribution); Lemma
B2 (the sign-birth channel: the owner's all-$p$ normalization idempotent law,
executed at $p=2$ on the discriminator, with the charge bookkeeping, the
D-identity, the sign law, and the typed bridge to the vertex); Theorem B3 (the
monodromy wall of the engine: transitive component monodromy forces total
W-death of every component --- certified on all transitive subgroups of
$S_2,\dots,S_5$ --- and the only remaining door is the involution-typed
discriminant split, the codimension-2 form of the committed quadric fold law);
Corollary B4 (the funnel bypass: the birth channel is not one of the six
funnel families and is not barred by any committed wall; plus the entry
checklist any candidate split must clear: grade-3 purity, FM-flatness, the
Weil phase circle, the $8m$ charge normalization). \textbf{NOT claimed:} any
decision of the vertex bit; any construction of a W-charged surface; any
exhaustiveness claim about all conceivable component families beyond the
certified monodromy constraints; any resolution of the campaign contract. This
strike \textbf{DETERMINES NO BIT}; it executes the owner's birth laws as a
typed construction channel at the vertex, certifies the channel's bookkeeping,
wall, and door, and arms the next strike (the involution-typed split on the
SAME carrier --- Round 18). The first unknown existence quantifier (a component
with nonzero Weil charge) is retained, exactly as the owner's Theorem B record
requires.

\section{The committed frame (what the funnel left)}

The campaign's certified reduction stack (cycles 2--4, rounds 15--16, all
receipts live): on the fixed test world --- the nonsplit $\mathbf{Q(i)}$-Weil
fourfold discriminator $\KA$ with $\operatorname{NS}(\KA)_{\mathbb Q} =
\mathbb Q\theta$, Weil classes $\WR, \WI$ (the parity-exact committed model:
$\int \theta^4 = 24$, $\int \WR^2 = \int \WI^2 = 8$, $\int \WR\WI = 0$,
$\WR\wedge\theta^j = 0$ for all $j \ge 1$ --- Lemma W0's class annihilation)
--- the vertex bit
\[
m^*(\KA,2,\WR) > 0
\]
is equivalent to: \emph{one W-charged surface exists} (a rational algebraic
codimension-2 class $Z$ with $\langle \WR, Z\rangle \ne 0$), equivalent to a
pure pair $(E_+,E_-)$ at some scale (c$_2$-difference $-m\WR$, Weil charges
differing by $8m$ --- the W6 table), by the committed funnel theorem and the
W2/W1 one-source equivalence (M3-normalized). The funnel certified that every
mechanism family of the committed classical arsenal (divisor intersections,
tensor/line-bundle transports, FM transforms, projector correspondences,
isogeny transport, Segre/pushforward on NS-sector cycles) generates $q = 0$
content only. On the discriminator the $(2,2)$-Hodge lattice is
$\mathbb Q\theta^2 \oplus \mathbb Q\WR \oplus \mathbb Q\WI$ (the committed
input stratification), so every surface class decomposes
\[
[D] \;=\; b\,\theta^2 \;+\; w\,\WR \;+\; w'\,\WI, \qquad
q(D) := \langle \WR, [D]\rangle = 8w \ \text{exactly}.
\]
The funnel killed the classical construction routes. What it did
\emph{not} touch --- because no committed theorem had ever executed it --- is
the owner's own birth arsenal. That execution is this round.

\section{The engine (the owner's committed birth laws)}

All four inputs are committed records of 2026-10-21--24 on this branch; the
strike adds no new mechanism, it runs these as engines:

\begin{itemize}
\item \textbf{Theorem A (all-$p$ normalization idempotents).} For any smooth
projective $X$ and any actual reduced pure codimension-$p$ scheme-theoretic
intersection $C = V(s_1,\dots,s_p) = \bigcup_{j=1}^m C_j$ with normalization
$\nu: \coprod_j \widetilde C_j \to X$: for every sign vector
$\sigma \in \{\pm1\}^m$ there is an algebraic involution $A_\sigma$ of the
normalization structure algebra acting by $\sigma_j$ on the $j$-th summand,
with finite idempotents $P_{\sigma,\pm} = (1 \pm A_\sigma)/2$, and the pushed
eigensheaf difference $\Xi_\sigma := [\nu_*(\operatorname{im}P_{\sigma,+})] -
[\nu_*(\operatorname{im}P_{\sigma,-})] \in K^0(X)$ satisfies the exact Chow
identity
\[
\ch^{\CH}_p(\Xi_\sigma) \;=\; \sum_{j=1}^m \sigma_j [C_j]
\quad\text{in } \CH^p(X)_{\mathbb Q},
\]
with $\ch_i(\Xi_\sigma) = 0$ for $i < p$ (the support filtration:
generic length one at each $C_j$).
\item \textbf{Theorem B (the sign law).} For a prescribed rational class $r$
coupling to the components by $a_j = \langle r, [C_j]\rangle$, the conductor
sign $\sigma^* = (\operatorname{sign} a_j)$ attains the maximum period
$\sum_j |a_j|$ over all $2^m$ sign vectors; it is nonzero iff some component
couples nontrivially. The record retains the first existence quantifier
explicitly.
\item \textbf{Theorem C (Betti rigidity).} For connected flat families of
pure codimension-$p$ cycles over a connected parameter space, the canonical
Betti class of each connected component family is constant; distinct component
classes arise only through discriminant crossings
(splitting/merging/monodromy).
\item \textbf{Theorem D (component charge conservation).} For same-degree
complete intersections, $\sum_j v_j = 0$ on primitive parts and
$\sum_j \sigma_j v_j = 2\sum_{\sigma_j = +1} v_j$; the component primitive
span has rank $\le m-1$.
\item \textbf{The quadric fold law.} On $Q = \{xy + z^2 = 1\}$ with
involution $A$, $P = (I+A)/2$: the born components of the principal
discriminant split satisfy $[L_+] = c_1^{\CH}(\operatorname{im}P)$ and
$[L_+] + [L_-] = 0$ while $[L_+] \ne 0$ --- the total is Chow-zero, the
signed component is the eigenbundle Chern class, in true integral Chow.
\end{itemize}

\section{Lemma B1: the W-dead total law}

\begin{lemma}[B1]
Let $C = V(s_1, s_2)$ be a pure codimension-2 complete intersection of
sections of line bundles $L_1, L_2$ on the discriminator with
$c_1(L_i) = k_i\theta$ (the NS-rank-one constraint). Then
\[
[C] \;=\; k_1 k_2\, \theta^2 \quad\text{in } \CH^2(\KA)_{\mathbb Q},
\qquad q(C) \;=\; \langle \WR, [C]\rangle \;=\; 0,
\]
and the annihilation is class-level: $\theta^2 \wedge \WR = 0$ (no surviving
monomials --- the committed W0). Consequently, for every reduced such $C$ with
components $C = \bigcup_j C_j$,
\[
\sum_{j=1}^m q_j \;=\; 0, \qquad q_j := \langle \WR, [C_j]\rangle.
\]
\end{lemma}

\begin{proof}
$[C] = c_1(L_1)\cdot c_1(L_2)$ by the intersection law; each
$c_1(L_i) \in \mathbb Z\theta$; the pairing $\langle \WR, \theta^2\rangle
= \int \WR \theta^2 = 0$ is the committed W0 annihilation (certified at
monomial level in the committed exterior model). Additivity of $q$ in the
cycle gives the component conservation.
\end{proof}

\emph{Certified:} the $81$-point $(k_1,k_2)$ grid with $-4 \le k_i \le 4$
($q(\text{total}) = 0$ everywhere) plus the monomial-level annihilation, in
the certificate script (Section~\ref{sec:verify}).

\emph{Content:} this is the codimension-2 form of the owner's whole-zero-locus
law (``the full Chow class of a section zero locus is fixed, independent of
the tuning''): the TOTAL of any NS-tuned intersection carries no Weil
information. Whatever Weil charge exists in the intersection lives only in
the \emph{signed distribution} among components --- the exact channel
Theorem A births.

\section{Lemma B2: the sign-birth channel and the typed bridge}

\begin{lemma}[B2]
Let $C = \bigcup_{j=1}^m C_j$ be any reduced pure codimension-2 complete
intersection on the discriminator (special sections allowed; the component
count $m \ge 1$ and the component geometry are unrestricted). Then:
\begin{enumerate}
\item[(i)] For every sign vector $\sigma \in \{\pm 1\}^m$, the owner's
Theorem-A eigensheaf difference $\Xi_\sigma \in K^0(\KA)$ is an actual
algebraic $K^0$ class with
\[
\ch^{\CH}_2(\Xi_\sigma) = \sum_j \sigma_j [C_j] \in \CH^2(\KA)_{\mathbb Q},
\qquad \ch_i(\Xi_\sigma) = 0 \ \ (i < 2),
\]
and its Weil charge is
\[
q(\Xi_\sigma) = \sum_j \sigma_j q_j = 2\!\!\sum_{\sigma_j = +1}\!\! q_j
\]
(the D-identity, holding because $\sum_j q_j = 0$ by B1).
\item[(ii)] The sign law (Theorem B with $r = \WR$): $\sigma^*_j =
\operatorname{sign}(q_j)$ attains
\[
q(\Xi_{\sigma^*}) \;=\; \sum_j |q_j| \;=\; \max_{\sigma} q(\Xi_\sigma).
\]
\item[(iii)] \textbf{The typed bridge.} The following are equivalent:
\begin{itemize}
\item[(a)] some reduced codimension-2 complete intersection on $\KA$ has a
component with $q_j \ne 0$;
\item[(b) ]a W-charged surface exists on $\KA$ (the component $C_j$ itself:
an actual algebraic surface with $\langle \WR, [C_j]\rangle = q_j \ne 0$);
\item[(c)] the vertex fires YES (the funnel's committed equivalence);
\item[(d)] a pure pair exists at some scale (the committed W2/W1
one-source reconstruction).
\end{itemize}
Moreover, under (a), the born class $\Xi_{\sigma^*}$ has charge
$\sum_j |q_j| \ge |q_{j^*}| > 0$: the engine births the charge the moment any
component carries it, with no further construction step.
\item[(iv)] \textbf{The one-source normalization.} If $[D] = b\theta^2 +
w\WR + w'\WI$ is any surface class, then $q(D) = 8w$ exactly; the pure pair
born by the one-source reconstruction at scale $m = |w|$ realizes the W6
charge split $-8m$ and the charged face $m\WR$ pairs to $8m$.
\end{enumerate}
\end{lemma}

\begin{proof}
(i) is Theorem A at $p = 2$ on $X = \KA$ (all hypotheses hold: $\KA$ smooth
projective, $C$ an actual reduced pure codimension-2 intersection); the charge
formula is the pairing of the Chow identity against $\WR$; the D-identity is
Theorem D's sign identity applied to the charge functional. (ii) is Theorem B
with $r = \WR$ (the periods $a_j = q_j$ are rational --- the pairing of a
rational Hodge class against algebraic classes). (iii): (a)
$\Rightarrow$ (b) is immediate; (b) $\Leftrightarrow$ (c) $\Leftrightarrow$
(d) are the committed funnel and W2/W1 equivalences (M3-normalized, both
directions attributed); (c) $\Rightarrow$ (a): if the vertex fires, a
W-charged surface $D$ exists; the committed Serre/Bertini component
cofinality theorem (2026-10-21) realizes every actual irreducible
codimension-2 algebraic $D$ as a multiplicity-one component of some
sufficiently ample complete intersection, giving (a). (iv): the pairing
computation on the committed stratification (Section 2), certified on the
$49$-point $(w,b)$ grid; the pair-face identities are the committed W6 table,
re-verified on born-charge instances in the script.
\end{proof}

\emph{Content:} the owner's birth law, run at the vertex, \emph{is} the
vertex's YES-channel: the equivalence (iii) is typed in both directions, and
the birth step (i)--(ii) converts any charged component into a charged
algebraic $K^0$ class by a finite explicit sign idempotent. What remains
unknown is exactly the owner's retained quantifier: whether any component of
any intersection on the discriminator carries nonzero charge. The next
theorem certifies where that quantifier can and cannot live.

\section{Theorem B3: the monodromy wall and the only door}

\begin{theorem}[B3]
Let $(C_t)_{t \in T}$ be a connected flat family of pure codimension-2
complete intersections on the discriminator over a connected base $T$, with
smooth locus $T^\circ$ and component-labeling $\pi_1(T^\circ)$-set. Then:
\begin{enumerate}
\item[(i)] \textbf{(Betti rigidity of the charge pattern).} The charge
function $q: \{\text{components over } T^\circ\} \to \mathbb Q$ is
$\pi_1(T^\circ)$-invariant: the analytic continuation of a component family
around any loop carries the same Weil charge.
\item[(ii)] \textbf{(The transitive kill).} If the component monodromy is
transitive, then all $q_j$ are equal, hence (by B1's conservation) all are
zero: \emph{transitive families birth no charge.}
\item[(iii)] \textbf{(The general member is W-dead).} A Bertini-general
member is irreducible ($m = 1$), and $q_1 = q(C) = 0$ by B1.
\item[(iv)] \textbf{(The only door).} By (i)--(iii), a nonzero component
charge can arise only where the component count changes --- at a discriminant
crossing --- and only with non-transitive local monodromy. The owner's quadric
fold law fixes the entry shape: an $m{:}1 \to 2$ split at the discriminant
with a local involution exchanging the two born components,
\[
[D_1] + [D_2] = [\text{generic fiber}] \ \ (\text{W-dead by B1}), \qquad
B := \tfrac{1}{2}\bigl([D_1] - [D_2]\bigr), \quad q(B) = q(D_1) = -q(D_2),
\]
the codimension-2 form of $[L_+] = c_1^{\CH}(\operatorname{im}P)$: the born
signed component carries the full charge of one face.
\item[(v)] \textbf{(The reducible-section stratum).} Components arising from
reducible divisors, $C = (D_1 + D_1') \cap (D_2 + D_2')$, are pairwise divisor
intersections --- funnel-dead, $q = 0$. The live candidate pools inside the
intersection channel are exactly: ($\alpha$) excess/reducible-intersection
components of intersections of irreducible divisors (not themselves expressible
as divisor intersections), and ($\beta$) the discriminant-split components of
(iv).
\end{enumerate}
\end{theorem}

\begin{proof}
(i) Theorem C: the Betti class of each connected component family is
constant on $T^\circ$; transport around a loop is an isomorphism of families,
preserving Betti classes, hence charges. (ii) A $\pi_1$-invariant function on
a transitive $\pi_1$-set is constant (single orbit); B1 forces the common
value $0$. Certified combinatorially: for every transitive subgroup
$G \le S_m$, $m = 2,\dots,5$ --- all $1 + 2 + 9 + 20 = 32$ of them, the
classical transitive-subgroup counts --- the $G$-invariant charge vectors form
exactly the constant line (exact rational nullspace computation per group).
(iii) Bertini irreducibility + B1. (iv) Theorem C (the component count and
component classes change only at discriminant crossings) plus the transitive
kill (ii) excludes transitive local degenerations; the involution-exchange
split is the committed quadric law's shape --- the two born components are
exchanged by the fold involution, the total is the W-dead generic class, and
the signed component is the eigenbranch difference. The charge identities are
B1 conservation plus the signed-component definition; certified on synthetic
split data ($q(D_1) = -q(D_2) = q(B) = 8w$ at scale $w$, with both
principal-type and NS-type totals). (v) Functoriality of the intersection
product on reducible divisors; the funnel's divisor-intersection family kill.
\end{proof}

\emph{Content:} this is the engine's own wall and its own door, certified
inside the birth channel --- not a wall of the ambient frame. The wall:
transitive monodromy and reducible-section geometry birth no charge. The
door: the involution-typed split. The wall sharpens the engine; the door aims
it.

\section{Corollary B4: the funnel bypass and the entry checklist}

\begin{corollary}[B4]
\begin{enumerate}
\item[(i)] \textbf{(Bypass.)} The birth channel --- Theorem-A sign idempotents
on intersection components, and the involution-typed discriminant splits of
B3(iv) --- is not one of the six funnel families: its outputs are signed
sums/differences of component classes, not products of divisor classes, not
tensor/FM/projector/isogeny/Segre images of NS-sector data. Any born class
with $q \ne 0$ lies outside the $q = 0$-only image of all six families: the
funnel does not bar this channel.
\item[(ii)] \textbf{(The entry checklist.)} Any candidate split of B3(iv)
whose born pair is to move the vertex must clear, in order:
\begin{itemize}
\item[\textbf{E1}] \emph{Grade purity:} $\ch_3(\Xi) = 0$ --- the
$c_1$-matching of the two components (the GRR grade-3 content
$\frac{1}{2}i_{j*}c_1(D_j)$ must cancel in the difference); grades $0,1$ die
automatically by the support filtration; grade 4 is the determined
c$_4$-constraint of W6.
\item[\textbf{E2}] \emph{FM-flatness:} $c_1(L_{\text{pair}}) = 0$ (the
committed W8 entry wall --- every future wall must be FM-closed).
\item[\textbf{E3}] \emph{The phase on the Weil circle:} the born pair's
monodromy phase $\psi = m\cdot\chi_W$ lands on the fixed rank-1 subtorus (the
committed W9 --- re-verified at the born scales $m = 1,\dots,4$ in the
script).
\item[\textbf{E4}] \emph{The charge normalization:} $q(D_1) = 8m$ at the
one-source scale $m$ (B2(iv)).
\end{itemize}
\end{enumerate}
\end{corollary}

\begin{proof}
(i) The six families' output classes are certified $q = 0$ (the funnel
theorem); a class with $q \ne 0$ is therefore not among their outputs; the
channel's outputs are the Theorem-A Chow identities and B3(iv) signed
components, which are neither products of divisor classes nor images of the
listed transports (the sign idempotent is a finite algebraic operation on the
normalization of a degeneration, not a correspondence of the listed types).
(ii) E1 is the purity requirement of the pair face (the W6 shared-$(r, c_1,
c_3)$ condition); E2--E3 are the committed W8/W9 entry walls, which every
pair-channel mechanism must respect (the dichotomy bit is duality-invariant
by W5, and the phase circle is the rank-1 residue of W9); E4 is B2(iv).
\end{proof}

\section{The armed next strike (Round 18: the split on the SAME carrier)}

The trajectory plan's Round 18 --- the integrability strike: construct
$(\bar\partial_E)^2 = 0$ on the SAME prescribed $E_\alpha$ from the GST
causal-fold laws alone; the committed Theorem B2 + descent law then give
$Z_\alpha$ with $\Cl(Z_\alpha) = \alpha$ exactly. The birth strike aims it:
the eigenbranch splitting of the fold's carrier IS an $m{:}1 \to 2$-type
birth exactly when the two branches' grade-2 contents carry opposite
Weil charges; B1--B4 convert the integrability strike's output into the
vertex's charge test directly --- E1--E4 are the checks the constructed
$Z_\alpha$ must pass, and the monodromy wall B3 certifies which degeneration
shapes can carry the charge at all (involution-typed, non-transitive). The
two-center determinantal formula $Z_\alpha$ (conditional only on
integrability) is the armed candidate: its branch geometry is
involution-typed by construction (the two centers), and the charge checklist
applies componentwise.

\section{The quantifier line (FS$_\forall$)}

The sign-birth law is input-typed and world-portable: on any world, the
per-input content of FS$_\forall$ (the bit $m^* > 0$) is the sign-detection
question on that world's component lattice --- $\sum_j |a_j| > 0$ with
$a_j$ the period pairings of the prescribed class against the components of
any sufficiently ample intersection (Theorem B with $r = \alpha$; the
Serre/Bertini cofinality makes the pool exhaustive over existing algebraic
components). Round 17 delivers: the executable per-world detection program
(the sign law + the one-source normalization), the discriminator's entry
shape (the only door), and the checklist. Round 19's transport sweep runs
this program world by world along the certified T-edges; each world
certified is a quantifier advance.

\section{Verification, labels, and scope discipline}
\label{sec:verify}

\begin{itemize}
\item \textbf{Certificate script:}
\texttt{research/experiments/task17a\_birth\_strike\_checks.py} --- bounded,
exact (Fraction/int coefficients; the parity-exact Weil construction; no
complex numbers, no floats, zero tolerance; every comparison an asserted
equality). Sections: the committed constants re-confirm; B1 (the 81-point
grid + the class-level annihilation); B2 (all $2^m$ sign patterns on eight
synthetic charge vectors, the D-identity, the sign-law maximization, the
49-point one-source normalization grid); B3 (all subgroups of $S_2$--$S_5$
enumerated by bounded closure, the classical transitive counts $1/2/9/20$
asserted, and for each of the 32 transitive groups the exact invariant
subspace is the constant line); B4 (the W6 pair table on four born-charge
instances including the detector face; the bypass statement); the split
bookkeeping (five synthetic involution splits, principal-type and NS-type
totals, the full-face charge identity); the W9 entry checklist (the FM
normalization $\widehat\Phi(\WR) = \widehat\WR$ exactly, the ghost census
$(0,0,m\WR,0,0)$ at $m = 1..4$, the $m$-linear phase). Runtime $< 2$ s.
\item \textbf{Labels:} B1--B4 extend the committed frame (the W/V labels
untouched; the owner's Theorems A--D cited, not relabeled).
\item \textbf{M3 attribution:} the bridge's equivalence content (iii) is the
committed funnel + W2/W1 (recorded as normalization); the new content is the
execution: B1's total law at the intersection level, B2's sign-birth
bookkeeping and typed bridge, B3's monodromy wall and door, B4's bypass and
checklist.
\item \textbf{Commitment:} research-only (this dossier, the record, the
INDEX entry, the script, the registry/STATUS movement); no \texttt{.lean}, no
\texttt{.yml} changes; no workflow fires expected (verified live at push).
\item \textbf{English-only and the banned-word gate:} run at commit time on
every file this round touches (ledger 076/081 discipline).
\end{itemize}

\end{document}
